\documentclass[10pt,a4paper]{amsart}
\usepackage[margin=19mm]{geometry}
\usepackage{amsmath,amssymb,mathtools}
\usepackage[scr=rsfs,cal=euler]{mathalfa}
\usepackage{xfrac}
\usepackage[unicode,hidelinks,bookmarks=false]{hyperref}
\newcommand{\ie}{i.e.\ }
\newcommand{\cf}{cf.\ }
\newcommand{\resp}{respectively\ }
\newcommand{\Subsection}{\subsection}
\providecommand{\nobreakdash}{\nobreak\hbox{-}}
\setlength{\emergencystretch}{2em}
\multlinegap0pt
\usepackage{bm}
\usepackage{bbm}
\usepackage{dsfont}
\usepackage{xspace}
\usepackage{cancel}
\usepackage{mathtools}
\usepackage{amsthm}
\makeatletter
\@ifundefined{theorem}{\newtheorem{theorem}{Theorem}}{}
\@ifundefined{proposition}{\newtheorem{proposition}[theorem]{Proposition}}{}
\makeatother
\DeclareMathOperator{\tr}{tr}
\newcommand{\R}{\mathbb{R}}
\newcommand{\tp}{{\scriptscriptstyle\mathsf{T}}}
\title[Generalized matrix nearness problems II]{Generalized matrix nearness problems II}
\author{Rongbiao Thomas Wang, Chi-Kwong Li,, Lek-Heng Lim}
\date{Source: arXiv:2605.30181v1 (CC BY 4.0)}
\begin{document}
\maketitle

\noindent\emph{Source and licence.} Generalized matrix nearness problems II. Version arXiv:2605.30181v1 (CC BY 4.0), distributed under the Creative Commons Attribution 4.0 International licence, as recorded in the arXiv licence metadata for this exact version and shown on the abstract page https://arxiv.org/abs/2605.30181v1. Full rights notice: licenses/arxiv-2605.30181-CC-BY-4.0.txt.\par\smallskip

\noindent\textit{Editorial scope.} Counted unit: the Proposition whose source label is \texttt{prop:complete} in the article (source0.tex lines 579--610, source label \texttt{prop:complete}), delivered together with the article's own complete original proof of it (source0.tex lines 612--692, one \texttt{proof} environment of 81 lines). No mathematics is written, repaired or re-derived: statement and proof are the article's own text line for line. Required context retained uncounted: none. Every definition, display and label of those lines is retained unchanged. Omitted from this package and not redistributed: 1--578, 611--611, 693--1426. Nothing inside a retained excerpt is omitted or altered: every retained body is delivered exactly as the article's lines are written, so no omission receipt of any kind accompanies this package. Citations: the retained excerpts carry no citation command at all, so no bibliography entry of the article is transcribed and no reference is redistributed. Reference adaptation: none was needed and none was made; every reference inside the delivered unit resolves inside it, so no unresolved reference or citation is produced. Printed slips kept literal: no printed word, symbol, hypothesis or number of the counted statement or of its proof was altered. Layout and class: the article's own class is \texttt{amsart} with packages including geometry, algorithm, algpseudocode, mathtools, pifont, enumitem, verbatim, tikz, filecontents, placeins, mlmodern, extarrows, graphicx, calc; the wrapper is the lane's pinned \texttt{amsart} editorial wrapper at 10pt on A4, which restates the theorem-like environments the retained text needs and adds the article's own macro definitions that the retained text needs (\texttt{\textbackslash R}, \texttt{\textbackslash tp}, \texttt{\textbackslash tr}), with extra wrapper packages (bm, bbm, dsfont, xspace, cancel, mathtools). The delivered numbering is the unit's own. Assets: no retained span contains a figure environment, a graphics-inclusion command, a PDF-page inclusion, a table or a TikZ picture; no image file is copied into this package, the package declares no asset dependency and no image file is redistributed with it. Licence: version arXiv:2605.30181v1 of this article is distributed under the Creative Commons Attribution 4.0 International licence, as recorded in the arXiv licence metadata for this exact version and shown on the abstract page https://arxiv.org/abs/2605.30181v1. A full rights notice is delivered at licenses/arxiv-2605.30181-CC-BY-4.0.txt. Characters: every retained excerpt is pure ASCII, so the delivered engine, XeLaTeX with its default Latin Modern text font, reports no missing character.
\begin{proposition}\label{prop:complete}
Let $\lVert\,\cdot\,\rVert$ be a norm on $\R^{m\times n}$ and $\lVert\,\cdot\,\rVert^*$ be the dual norm. For any $S \in \R^{m \times n}$, we write $\varphi_S : \R^{m \times n} \to \R$, $X \mapsto \tr(S^\tp X)$, for the linear functional defined by $S$. Let $A_{12}\in \R^{s \times (n-t)}$, $A_{21}\in \R^{(m-s) \times t}$, $A_{22}\in \R^{(m-s) \times (n-t)}$, and $Y_* \in \R^{s \times t}$. 
Then
\[
\biggl\lVert \begin{bmatrix}
    Y_* & A_{12} \\
A_{21} & A_{22}
\end{bmatrix} \biggr\rVert \le \biggl\lVert \begin{bmatrix}
    Y & A_{12} \\
A_{21} & A_{22}
\end{bmatrix} \biggr\rVert \quad \text{for all } Y \in \R^{s \times t}
\]  if and only if there exists 
\begin{equation}\label{eq:R_form}
    S= \begin{bmatrix}
    0 & S_{12} \\
S_{21} & S_{22}
\end{bmatrix} \in \R^{m \times n}
\end{equation}
with $S_{12}\in \R^{s \times (n-t)}$, $S_{21}\in \R^{(m-s) \times t}$, $S_{22}\in \R^{(m-s) \times (n-t)}$ such that $\lVert \varphi_S \rVert^*=1$ and 
\begin{equation}\label{eq:completion_thm}
    \biggl\lVert \begin{bmatrix}
        Y_* & A_{12} \\
    A_{21} & A_{22}
    \end{bmatrix} \biggr\rVert = \tr\biggl(S^\tp \begin{bmatrix}
        Y_* & A_{12} \\
    A_{21} & A_{22}
    \end{bmatrix} \biggr) = \tr \biggl( S^\tp  \begin{bmatrix}
        0 & A_{12} \\
    A_{21} & A_{22}
    \end{bmatrix}  \biggr).
\end{equation}
\end{proposition}

\begin{proof}
Let $\beta \coloneqq \bigl \lVert \bigl[\begin{smallmatrix}
    Y_* & A_{12} \\
A_{21} & A_{22}
\end{smallmatrix}\bigr]  \bigr\rVert$. By its definition, $\beta \le \bigl \lVert \bigl[\begin{smallmatrix}
    Y & A_{12} \\
A_{21} & A_{22}
\end{smallmatrix}\bigr]  \bigr\rVert$ for all $Y \in \R^{s \times t}$. Consider the subspace
\[
\mathbb{V} = \biggl\{ \begin{bmatrix}
    Y &zA_{12}\\
    zA_{21} &zA_{22}
\end{bmatrix} \in \R^{m \times n} : Y \in \R^{s\times t},\, z \in \R  \biggr\}
\]
and the linear functional $\varphi : \mathbb{V} \to \mathbb{R}$ defined by
\[
\varphi\biggl(\begin{bmatrix} 
    Y &zA_{12}\\
    zA_{21} &zA_{22}
\end{bmatrix}\biggr) = \beta z
\]
for all $Y \in \R^{s\times t}$ and $z \in \R$. 
We must also have $\lVert \varphi\rVert^* \le 1$ as
\[
\biggl\lVert \begin{bmatrix}
    Y &zA_{12}\\
    zA_{21} &zA_{22}
\end{bmatrix} \biggr\rVert = \lvert z \rvert \biggl\lVert \begin{bmatrix}
    Y/z &A_{12}\\
    A_{21} & A_{22}
\end{bmatrix} \biggr\rVert \ge \lvert  \beta z\rvert =  \biggl\lvert \varphi \biggl( \begin{bmatrix} 
    Y &zA_{12}\\
    zA_{21} &zA_{22}
\end{bmatrix} \biggr) \biggr\rvert
\]
for all $Y\in \R^{s \times t}$  and $0 \ne z \in \R$.
In fact, $\lVert  \varphi\rVert^* = 1$ since $\beta =  \varphi\bigl(  \bigl[\begin{smallmatrix}
    Y_* & A_{12} \\
A_{21} & A_{22}
\end{smallmatrix}\bigr] \bigr) \le \lVert \varphi \rVert^* \beta$. By Hahn--Banach Theorem, $\varphi$ can be extended to a linear functional on $\R^{m\times n}$ with the same norm. By Riesz representation, any such linear functional must take the form $\varphi_S$ for some $S  \in \R^{m\times n}$.  So  $\lVert \varphi_S\rVert^* = 1$ and $\varphi_S \bigr\rvert_{\mathbb{V}} =\varphi$, i.e.,
\[
\varphi_S\biggl( \begin{bmatrix}
    Y &zA_{12}\\
    zA_{21} &zA_{22}
\end{bmatrix} \biggr) =
\tr\biggl( S^\tp \begin{bmatrix}
    Y &zA_{12}\\
    zA_{21} &zA_{22}
\end{bmatrix} \biggr)=\beta z
\]
for all $Y\in \R^{s \times t}$ and $z \in \R$. Since 
$
\tr \bigl( S^\tp 
\bigl[\begin{smallmatrix}
    Y &0\\
    0 &0
\end{smallmatrix}\bigr] \bigr) = 0
$
for all $Y \in \R^{s \times t}$, the matrix $S$ must be of the form \eqref{eq:R_form} and satisfies \eqref{eq:completion_thm}.

Now, suppose that there exists $S\in \R^{m \times n}$ of the form \eqref{eq:R_form} such that $\lVert \varphi_S \rVert^*=1$ and satisfies \eqref{eq:completion_thm}. Then for any $Y\in \R^{s \times t}$,
\begin{align*}
    \biggl\lVert  \begin{bmatrix}
        Y_* & A_{12} \\
    A_{21} & A_{22}
    \end{bmatrix}  \biggr\rVert &= \biggl\lvert \tr\biggl(S^\tp  \begin{bmatrix}
        Y_* & A_{12} \\
    A_{21} & A_{22}
    \end{bmatrix} \biggr) \biggr\rvert \\
    &= \biggl\lvert \tr\biggl(S^\tp  \begin{bmatrix}
        Y & A_{12} \\
    A_{21} & A_{22}
    \end{bmatrix} \biggr) \biggr\rvert\le \lVert S\rVert^* \biggl\lVert  \begin{bmatrix}
        Y & A_{12} \\
    A_{21} & A_{22}
    \end{bmatrix}  \biggr\rVert = \biggl\lVert  \begin{bmatrix}
        Y & A_{12} \\
    A_{21} & A_{22}
    \end{bmatrix}  \biggr\rVert. \qedhere
\end{align*}
\end{proof}

\end{document}
